
%----------------------------------------------------------------------------------------
%	DOCUMENT DEFINITIONS
%----------------------------------------------------------------------------------------

\documentclass[a4paper]{resume}

\usepackage[left=0.75in,top=0.6in,right=0.75in,bottom=0.6in]{geometry}
\usepackage[T1]{fontenc}
\usepackage[utf8]{inputenc}	
\usepackage[english]{babel}
\usepackage{microtype}
\usepackage{lmodern}
\usepackage{enumitem}
\usepackage{multicol}
\usepackage{booktabs}
\usepackage{hyperref}

\hypersetup{
    colorlinks=true,
    linkcolor=blue,
    filecolor=magenta,      
    urlcolor=black,
    pdftitle={Pedro Henrique Rosso - Resume},
    }

\urlstyle{same}

\newcommand{\tab}[1]{\hspace{.2667\textwidth}\rlap{#1}}
\newcommand{\itab}[1]{\hspace{0em}\rlap{#1}}

\name{Pedro Henrique Di Francia Rosso}
\address{Boaventura do Amaral, 1360 (Apto 106), Centro, Campinas - SP, Brasil, CEP 13015-192}
\address{+55 (48) 99956-8129 \\ {\tt pedro.rosso@ic.unicamp.br | pedrohrosso@gmail.com}}

\begin{document} 

%----------------------------------------------------------------------------------------
%	EDUCATION SECTION
%----------------------------------------------------------------------------------------

\begin{rSection}{Education}
{\bf Computer Engineering} \hfill {\em 2014--2018} \\
{\sc Federal University of Santa Catarina (UFSC)} \hfill {\em Araranguá-SC, Brasil} \\
Participated in programming marathons, course monitor and scientific iniciation fellow. 
As undergrad final project, I studied image re-coloration for dichromats.
Best student of Engineering classes (Computer and Energy). GPA 9.04/10.0.

{\bf Master in Computer Science} \hfill {\em 2019--2021} \\
{\sc Federal University of ABC (UFABC)} \hfill {\em Santo André-SP, Brasil} \\
Grade A (maximum) in every class. I've presented the thesis: ``OCFTL: {\it an MPI implementation-independent
fault tolerance library for task-based applications}'' in July of 2021. Supervisioned by Emilio Francesquini.

{\bf PhD in Computer Science} \hfill {\em 2021--Current} \\
{\sc Universidade Estadual de Campinas (Unicamp)} \hfill {\em Campinas-SP, Brasil} \\
Currently enrolled as PhD student at UNICAMP, researching FPGAs programming with OpenMP in clusters. 
Supervisioned by Guido Araújo (UNICAMP) and Emilio Francesquini (UFABC).
\end{rSection}

%----------------------------------------------------------------------------------------
%	SKILLS SECTION
%----------------------------------------------------------------------------------------

\begin{rSection}{Awards and Distinctions}

%{\bf Honorable mention} $\diamond$ {UFSC} \hfill{\em 2015--2018} \\
%In the last 4 semesters of undergrad, received the Honorable mention
%by UFSC by obtaining the average semestral coeff above 9.0. From 6$^o$ to the 9$^o$ semester.

{\bf Best Academic Performance} $\diamond$ {UFSC} \hfill{\em 2018} \\
In February of 2019, received a medal from Federal University of Santa Catarina for getting the best 
performance in the Computer Engineering course -- 2018.

{\bf Best Student of Engineering Courses} $\diamond$ {CREA-SC} \hfill{\em 2018} \\
In February of 2019, received a distinction of the Regional Council of Engineering and Agronomy
for getting the best academic performance between the Computer and Energy Engineering courses -- 2018.

{\bf Paper Awards} $\diamond$ {ERAD-SP} \hfill{\em 2020 e 2021} \\
In ERAD-SP edition of 2020 and 2021, received the best graduate paper award.

{\bf Honorable Mention} $\diamond$ {WSCAD-CTD} \hfill{\em 2021} \\
In the WSCAD-CTD edition of 2021, received an honorable mention for the master thesis finishing among the three finalist thesis in the competition. 

\end{rSection}

\begin{rSection}{Research Interests}
Following, a non-exaustive list of my research interests:

\setlength\multicolsep{5pt}
\begin{multicols}{2}
\begin{itemize}[noitemsep]
  \item High-Performance Computing;
  \item Cluster programming
  \item FPGAs
  \item Fault Tolerance;
  \item Parallel and Distributed Computing;
\end{itemize}
\end{multicols}

\end{rSection}

%----------------------------------------------------------------------------------------
%	EXPERIENCE SECTION
%----------------------------------------------------------------------------------------

\begin{rSection}{Experience}

%{\bf Data Structure class monitor} \hfill{\em 2016} \\
%{\sc Universidade Federal de Santa Catarina (UFSC)} \hfill {\em Araranguá-SC, Brasil} \\
%In the 2 semesters of 2016 I monitored the students of the Data Structures discipline, 
%where did content reviews, exercises and tests. Approved with maximum grade.

{\bf Scientific Iniciation fellow} \hfill{\em 2017--2018} \\
{\sc Universidade Federal de Santa Catarina (UFSC)} \hfill {\em Araranguá-SC, Brasil} \\
During the 2nd semester of 2017 and 1st semester of 2018 I participated in a Scientific Initiation project, where a wireless electromyography application was developed. Funded by CNPq.

{\bf Master fellow} \hfill{\em 2019--2021} \\
{\sc Universidade Federal do ABC (UFABC)} \hfill {\em Santo André-SP, Brasil} \\
From the year 2019 to 2021, conducted my master's research focusing on fault tolerance for MPI. Master's work funded by FUNCAMP. During the same period, I was also a participant in the OmpCluster project under the coordination of prof. Guido Araújo (UNICAMP).

{\bf PhD fellow} \hfill{\em 2021--Current} \\
{\sc Universidade Estadual de Campinas (Unicamp)} \hfill {\em Santo André-SP, Brasil} \\
Since the year 2021 I have been doing a PhD focused on integrating FPGA programming into cluster programming using OpenMP. PhD work funded by FAPESP (2021/09355-2). Participating in the OmpCluster project under the coordination of prof. Guido Araújo (UNICAMP).

{\bf Research Intern} \hfill{\em Jan/2023--July/2023} \\
{\sc Advanced Micro Devices (AMD)} \hfill {\em Dublin, Ireland} \\
Integration of FPGA programmin with OpenMP with ACCL - Alveo Collective Communication Library. Colaboration with the Resarch Labs of AMD/Xilinx Ireland.
\end{rSection}
%----------------------------------------------------------------------------------------
%	SKILLS SECTION
%----------------------------------------------------------------------------------------

\begin{rSection}{Skills and Language Proficiency}
{\bf Programming Languages and Frameworks}

\begin{tabular}{ll}
  More experience & C, C++ \\
  Intermediate experience & Java, JavaScript, Node.js\\
  Familiarity & Assembly, C\#, Bash, CUDA, Javascript, Robotic \\
\end{tabular}

%{\bf Tools}
%\begin{itemize}[noitemsep]
%  \item Linux Operating System and utilities ({\tt grep}, {\tt
%    bash}, {\tt make}, etc);
%  \item OpenMP e MPI;
%  \item Source control ({\tt git}) and plataforms (Github, Gitlab,
%    etc);
%  \item Embedded Computing (project, prototyping e software).
%  \item MatLab/Octave.
%  \item Debuggers ({\tt gdb}) e perfiladores ({\tt perf}, {\tt strace}, {\tt
%    ltrace});
%  \item LLVM compiler.
%  \item {\tt gcc} and {\tt clang} C/C++ compilers;
%\end{itemize}

{\bf Language Proficiency}
\begin{itemize}[noitemsep]
  \item {\bf Portuguese} (Native) -- understand well; speak well; write well.
  \item {\bf English} (Advanced) -- understand well; speak well; write well.
\end{itemize}
\end{rSection}

\begin{rSection}{Tutorials}

{\bf 12$^o$ INFIERI School} $\diamond$ São Paulo \hfill {\em 2023} \\
Author. Hand-on Labs on FPGA-based acceleration of scientific applications using OpenMP Cluster - \textit{School of Intelligent Signal Processing for Frontier Research and Industry }.

\end{rSection}

\begin{rSection}{Event Participations}

{\bf 12$^o$ CARLA} $\diamond$ México, Virtual \hfill {\em 2021} \\
Author~\cite{carla2021}. Latin America High Performance Computing Conference.

{\bf 12$^a$ ERAD-SP} $\diamond$ UFABC/USP, Virtual \hfill {\em 2021} \\
Author~\cite{erad2021}. Escola Regional de Alto Desempenho do Estado de São Paulo.

{\bf 11$^a$ ERAD-SP} $\diamond$ UNESP/Mackenzie/USP, Virtual \hfill {\em 2020} \\
Author~\cite{erad2020}. Escola Regional de Alto Desempenho do Estado de São Paulo.

\end{rSection}

%\newpage
%\begin{rSection}{Links}
%  \begin{tabular}{ll}
%    Homepage       & \url{https://pedroohr.github.io} \\
%    Github         & \url{https://github.com/pedroohr} \\
%    Gitlab         & \url{https://gitlab.com/phrosso} \\
%    LinkedIn       & \url{https://twitter.com/pedrohrosso} \\
%    Google Scholar & \url{https://scholar.google.com/citations?user=rtONezgAAAAJ} \\
%    Lattes         & \url{http://lattes.cnpq.br/5343894554617060} \\
%  \end{tabular}
%\end{rSection}
\nocite{*}
\begin{rSection}{Selected Publications}
\bibliographystyle{abbrv}
\renewcommand{\section}[2]{}
\bibliography{./reference}
\end{rSection}

\vfill \hfill {\em Updated in: \today.}

\end{document}